%% FORMLÆRE: LÆREBOK — Kapittel 6–9 + Appendiks
%% Denne fila vert \input{} i hovuddokumentet.
%% Iver Raknes Finne, Institutt for design, AHO
%% Utkast 02 — 2026-03-30
%%
%% Avhengig av teorem-miljø, kommandoar og pakkar definerte i hovudfila.
%%
%% ENDRINGAR frå utkast_01:
%%  - Kap 7 og 8 er bytte: Materialsignaturar (empirisk) kjem no FØR
%%    Variasjonsprinsipp.
%%  - Spreiingsrelasjonen er nedgradert frå Teorem til Proposisjon.
%%  - Kompileringsfeil (missing \item) fiksa.
%%  - Læringsmål lagt til i kvart kapittel.
%%  - Redundans kutta (STOLAR-data berre her, posisjonsinvarians-døme
%%    konsolidert, Kubler berre éin gong, MI-utrekning komprimert).
%%  - Plast lagt til i materialtabell.
%%  - Ikkje-vestlege døme lagt til.
%%  - Mahogni-dateringar retta (1770, ikkje 1750). Art Nouveau: 1914.
%%  - Thonet nr. 14, Breuer B3, Eames DCW/Lounge lagt til.
%%  - Øvingshint kutta til maks to setningar.
%%  - Tabellar med \footnotesize der naudsynt.
%%  - Xenobotar-formulering presisert.
%%  - Boltzmann-antaknad eksplisitt markert i spreiingsbeviset.
%%  - Overflødig appendiks fjerna (berre éin korrespondansetabell).


%% =====================================================================
%%  KAPITTEL 6: SUBSTRAT-UAVHENGIGHEIT OG EMERGENS
%% =====================================================================
\chapter{Substrat-uavhengigheit og emergens}
\label{chap:substrat}

\begin{quote}\itshape
Etter dette kapittelet skal lesaren kunne:
\begin{itemize}[nosep]
\item Formulere postulatet om substrat-uavhengig navigasjon og kva som ville falsifisere det.
\item Bevise posisjonsinvarians: posisjonen i formrommet er uavhengig av produksjonsprosessen.
\item Forklare emergent navigasjon og kvifor det er ein konjektur, ikkje eit teorem.
\item Samanlikne navigatorar på tvers av substrat i eit felles formrom.
\end{itemize}
\end{quote}

\begin{quote}
\itshape
Materialet varierer. Logikken er identisk.
\end{quote}

\section{Navigasjon er ein eigenskap ved organisering}

Ein termostat oppfyller alle tre vilkåra frå navigatordefinsjonen
(kap.~5). Det same gjer ein gradient-descent-algoritme, ein koloni
av celler som byggjer eit organ, og ein handverkar som dreier ein
stolbein. Substratet --- metall, silisium, organisk vev,
menneskeleg medvit --- varierer radikalt. Den formelle strukturen
gjer det ikkje. Navigasjon er ein eigenskap ved \emph{korleis} eit
system er organisert, ikkje ved \emph{kva} det er laga av.


\section{Postulat: Substrat-uavhengig navigasjon}

\begin{postmark}[Substrat-uavhengig navigasjon]\label{post:L.7.1}
Vilkåra (a), (b) og (c) i definisjonen av ein navigator legg ingen
restriksjon på det fysiske substratet. Eitkvart fysisk system ---
biologisk, algoritmisk, mekanisk eller distribuert --- som:
\begin{enumerate}[label=(\roman*),nosep]
  \item har tilstandar (posisjonar i eit konfigurasjonsrom),
  \item kan registrere avstand frå ein målkonfigurasjon, og
  \item kan justere tilstanden sin som respons,
\end{enumerate}
kan oppfylle navigatordefinisjonens vilkår.
Navigasjon er ein eigenskap ved den \emph{funksjonelle
organiseringa} til systemet, ikkje ved materialet det er laga av.
\end{postmark}

\paragraph{Falsifiseringsvilkår.}
Postulatet fell dersom det kan visast at minst eitt av vilkåra
(a)--(c) \emph{nødvendigvis} krev eit spesifikt fysisk substrat.
Meir presist: dersom det finst formgjevingsprosessar der ein
menneskeleg navigator konsekvent produserer former som ingen
algoritmisk navigator kan produsere, \emph{ikkje} på grunn av
reknekapasitet eller sensorisk oppløysing, men fordi sjølve
navigasjonen krev organisk substrat, då fell postulatet.

Skilnaden mellom ``endå ikkje realisert'' og ``prinsipielt
umogleg'' er skilnaden mellom ingeniørarbeid og falsifisering.


\section{Posisjonsinvarians}

\begin{satmark}[Posisjonsinvarians]\label{thm:L.7.2}
Posisjonen $\bm{x} \in \MC$ til eit realisert objekt er bestemt
av dei målbare eigenskapane til objektet, ikkje av prosessen som
produserte det. To objekt med identiske eigenskapsvektorar
okkuperer same punkt i $\MC$, uansett om dei vart laga for hand,
av maskin, av algoritme, eller av biologisk vekst.
\end{satmark}

\begin{proof}[Bevis, steg for steg]
\begin{enumerate}
  \item \textbf{Utgangspunkt.}\;
    Eigenskapsvektoren $\bm{x}(a) = (x_1(a), \ldots, x_n(a))$ er
    definert utelukkande av dei målbare geometriske eigenskapane
    til objektet~$a$.
  \item \textbf{Nøkkelobservasjon.}\;
    Desse eigenskapane er \emph{intrinsiske} for objektet. Dei er
    ikkje funksjonar av produksjonshistoria. Eit stykke tre har
    same kurvatur anten det vart forma med kniv eller med
    CNC-fres.
  \item \textbf{Konklusjon.}\;
    La $a$ vere eit handlaga objekt og $b$ eit maskinlaga objekt.
    Dersom $x_i(a) = x_i(b)$ for alle $i = 1, \ldots, n$, so er
    $\bm{x}(a) = \bm{x}(b)$, og dei to objekta okkuperer same
    punkt i~$\MC$. \qedhere
\end{enumerate}
\end{proof}

\paragraph{Kva teoremet ikkje seier.}
Posisjonsinvarians betyr \emph{ikkje} at produksjonsprosessen er
irrelevant. Prosessen avgjer kva \emph{region} av $\MC$ som er
tilgjengeleg --- substratet avgrensar trajektorien. Men for dei
punkta som \emph{er} tilgjengelege, er posisjonen
sjølv substrat-invariant.


\section{Emergent navigasjon}

\begin{konmark}[Emergent navigasjon]\label{conj:L.7.3}
La $\Nav^1, \Nav^2, \ldots, \Nav^S$ vere ei samling
sub-navigatorar. Det kollektive systemet kan realisere eit
navigator-trippel $(G, \mu, \alpha)$ der $G$, $\mu$ og $\alpha$
\emph{ikkje} er funksjonar av nokon einskild sub-navigators
trippel, men oppstår frå interaksjonsdynamikken.
\end{konmark}

\paragraph{Den formelle strukturen.}
Kvar sub-navigator $\Nav^j$ genererer ein straum
$\dot{\bm{x}}^j = \alpha_j(\bm{x}^j)$. Når dei interagerer vert
den kollektive dynamikken:
\[
  \dot{\bm{x}}^j = \alpha_j(\bm{x}^j)
  + \sum_{l \neq j} \phi_{jl}(\bm{x}^j, \bm{x}^l)
\]
der $\phi_{jl}$ kodar interaksjonen mellom sub-navigatorane $j$
og~$l$. Til dømes: $\phi_{jl}(\bm{x}^j, \bm{x}^l)
= \kappa (\bm{x}^l - \bm{x}^j)$, ein enkel attraksjonskopling.

Det kollektive systemet har ein makroskopisk navigator
$(G, \mu, \alpha)$ dersom det finst ei grovkorning (ein projeksjon
frå mikro- til makrotilstandar)
$\Pi : (\MC)^S \to \MC$ slik at den projiserte dynamikken
$\dot{\bm{y}} = \Pi_* \dot{\bm{X}}$ oppfyller
navigatordefinisjonens tre vilkår --- og trippelet ikkje kan
dekomponerast som ein funksjon av nokon einskild
$(G_j, \mu_j, \alpha_j)$.

\paragraph{Kvifor er dette ein konjektur og ikkje eit teorem?}
Vi påstår at slik emergens \emph{finst}, men vi har ikkje
presisert vilkåra på $\phi_{jl}$ som garanterer det.

\paragraph{Empirisk grunngjeving.}
Syntetiske organismar --- xenobotar --- sette saman av
froskeembryoceller, navigerer mot koherente morfologiske
konfigurasjonar. Cellene har individuell evolusjonshistorie (dei
er celler frå \emph{Xenopus laevis}), men den kollektive
konfigurasjonen deira har det ikkje. Navigasjonskapasiteten
oppstår frå intercellulær bioelektrisk signaldynamikk.


\section{Designhistorisk døme: Konvergens på tvers av substrat}

I 1917 skildra D'Arcy Thompson korleis biologiske
beinkonstruksjonar konvergerer mot former som minimerer
materialebruk under mekanisk last. Hundre år seinare produserer
topologisk optimering nesten identiske former under same fysiske
gradientar.

\medskip
\begin{center}
{\footnotesize
\begin{tabular}{@{}lp{30mm}p{30mm}@{}}
\toprule
& \textsc{Biologisk bein} & \textsc{Topologisk opt.} \\
\midrule
Substrat & Organisk vev & Metall/polymer \\
Navigator & Cellulær signalering & Algoritme (SIMP/ESO) \\
Seleksjonstrykk & Mekanisk last + metabolsk kostnad
  & Stivleik + massebudsjett \\
Resulterande form & Trabekulært nettverk & Lattice-struktur \\
Konvergens & \multicolumn{2}{c}{Same topologi under same gradientar} \\
\bottomrule
\end{tabular}
}
\end{center}

\medskip
Mingdynasti-møblar illustrerer same prinsipp i ein annan skala:
presise samanføyingar (mortise-and-tenon utan lim) i hardtre
konvergerer mot same geometriske løysingar som moderne CNC-fresa
tresamband --- trass i radikalt ulike navigatorar (handverkar vs.\
algoritme) og seks hundre års avstand.


\section{Gjennomarbeidd døme: Tre navigatorar, éin stol}

Tre stolformer produserte under same funksjonelle krav (bere
80~kg, sitjehøgd 43--47~cm, stabil på plant golv) og same
estetiske preferanse (visuell lettheit):

\paragraph{Navigator A: Handverkaren.}
Ein snikkar med 20 års erfaring i bøk. Resultatet er ein stol med
lett kurva bakstykke, runde stolbeinsnitt, og ein kompaktheit
rundt 0{,}44.

\paragraph{Navigator B: Algoritmen.}
Ein topologisk optimeringsprosess. Resultatet er ein stol med
organisk, forgreina struktur og ein kompaktheit rundt 0{,}41.

\paragraph{Navigator C: Evolusjonen.}
Tusen år med stoltradisjon. Resultatet er ein vernakular stoltype
med rette linjer og ein kompaktheit rundt 0{,}46.

\paragraph{Samanlikning.}
\emph{Merk: kompaktheitsverdiane er fiktive, valde for å illustrere
prinsippet. Empiriske verdiar kjem i kapittel~7.}

\begin{center}
{\footnotesize
\begin{tabular}{@{}lccc@{}}
\toprule
& \textsc{Handverkar} & \textsc{Algoritme} & \textsc{Evolusjon} \\
\midrule
Kompaktheit & 0{,}44 & 0{,}41 & 0{,}46 \\
Symmetri & 0{,}12 & 0{,}09 & 0{,}14 \\
Kurvatur & 0{,}10 & 0{,}13 & 0{,}06 \\
\midrule
Substrat & Bøk & Polymer/metall & Bøk \\
Trajektorie & Lokal, inkr. & Sprangvis & Blind, kumulativ \\
\bottomrule
\end{tabular}
}
\end{center}

\medskip
Alle tre ligg innanfor same \emph{region} av formrommet fordi
seleksjonstrykka er dei same. Men dei okkuperer \emph{ulike punkt}
fordi substrata (og dermed trajektoriane) er ulike.
Substratet formar stigen, ikkje målet.


\section{Øvingar}

\begin{exercise}[Nivå 1]\label{ex:6.1}
Ein termostat, ein gradient-descent-algoritme, og ein bie som
byggjer ein vokskake: identifiser $(G, \mu, \alpha)$ for kvar av
dei tre navigatorane. Kva er likt? Kva er ulikt?
\end{exercise}

\begin{exercise}[Nivå 1]\label{ex:6.2}
Forklar med eigne ord kvifor posisjonsinvarians \emph{ikkje} betyr
at produksjonsprosessen er irrelevant.
\end{exercise}

\begin{exercise}[Nivå 2]\label{ex:6.3}
Eit designteam på fem personar utviklar ein ny kontorstol. Bruk
konjekturen om emergent navigasjon til å modellere teamet som eit
kollektiv av sub-navigatorar. Kva er $\phi_{jl}$ i dette
tilfellet?
\end{exercise}

\begin{exercise}[Nivå 2]\label{ex:6.4}
Topologisk optimering av ein bjelke under bøyelast og biologisk
beinvekst under same lasttilfelle produserer liknande strukturar.
Kva eigenskapar ($x_1, \ldots, x_n$) ville du bruke for å
samanlikne dei i eit felles formrom?
\end{exercise}

\begin{exercise}[Nivå 3]\label{ex:6.5}
Design eit eksperiment som ville avgjere om ein spesifikk
formgjevingsoppgåve \emph{prinsipielt} krev organisk substrat.
Korleis ville du skilje mellom ``endå ikkje realisert'' og
``prinsipielt umogleg''?
\end{exercise}

\begin{exercise}[Nivå 3]\label{ex:6.6}
Finn eitt publisert døme (frå biologi, robotikk, eller
designforsking) der eit kollektivt system viser
navigasjonskompetanse som ikkje kan reduserast til nokon einskild
komponent. Spesifiser $(G, \mu, \alpha)$ og argumenter for kvifor
dei er emergente.
\end{exercise}


%% =====================================================================
%%  KAPITTEL 7: EMPIRISK — MATERIALSIGNATURAR
%%  (Tidlegare kap. 8 — flytta hit per PED/W-tilråding)
%% =====================================================================
\chapter{Empirisk: Materialsignaturar}
\label{chap:empirisk}

\begin{quote}\itshape
Etter dette kapittelet skal lesaren kunne:
\begin{itemize}[nosep]
\item Formulere postulatet om geometrisk materialsignatur og kva som ville falsifisere det.
\item Bevise informasjonsteoremet: materialet ber meir informasjon om form enn funksjonen.
\item Rekne ut gjensidig informasjon for eit enkelt datasett.
\item Tolke dei empiriske MI-verdiane frå STOLAR-datasettet.
\end{itemize}
\end{quote}

\begin{quote}
\itshape
Materialet er ikkje nøytralt. Det kanaliserer forma.
\end{quote}

\section{Frå teori til data}

Kapitla~1--6 bygde opp eit deduktivt rammeverk. Dette kapittelet
testar rammeverket mot data. Det presenterer det sterkaste
enkelttestresultatet og viser korleis det passar inn.


\section{Postulat om geometrisk signatur}

\begin{postmark}[Geometrisk materialsignatur]\label{post:L.6.1}
For kvart materiale $m$ brukt i ein funksjonell klasse $C$ er den
betinga fordelinga av geometriske eigenskapar gjeve materiale
ikkje-uniform:
\[
  P(\bm{x} \mid m) \;\neq\; \mathrm{Uniform}(\MC)
  \qquad \forall\; m.
\]
Kvart materiale induserer ein karakteristisk geometrisk
\emph{profil}: ein region i formrommet som det favoriserer.
\end{postmark}

\paragraph{Falsifiseringsvilkår.}
Postulatet fell om same geometriske distribusjon kan observerast
uavhengig av materiale, innanfor same funksjonelle klasse.

\paragraph{Kva betyr dette?}
Materialet \emph{kanaliserer} forma. Mahogni trekkjer mot andre
geometriar enn stål, ikkje berre fordi designaren vel det, men
fordi dei innbygde eigenskapane til materialet --- brotstyrke,
bøyelegheit, bearbeidingskarakter --- favoriserer visse
konfigurasjonar og motset seg andre.


\section{Informasjonsteoremet}

\begin{satmark}[Informasjonsteoremet]\label{thm:L.6.2}
Innanfor ein funksjonell klasse $C$ med minst to materialkategoriar
ber materialet strengt meir gjensidig informasjon med realisert
geometri enn funksjonen gjer:
\[
  \MI(G;\, M) \;>\; \MI(G;\, \mathrm{Func}) \;=\; 0.
\]
\end{satmark}

\begin{proof}[Bevis, steg for steg]
\paragraph{Del~1: $\MI(G;\, \mathrm{Func}) = 0$.}%
\begin{enumerate}
  \item \textbf{Premiss.}\;
    Innanfor ein funksjonell klasse $C$ deler alle objekt same
    funksjon per definisjon (kap.~2).
  \item \textbf{Konsekvens.}\;
    Variabelen $\mathrm{Func}$ er \emph{konstant} innanfor $C$.
    Ein konstant variabel har null entropi:
    $\Ent(\mathrm{Func}) = 0$.
  \item \textbf{Gjensidig informasjon.}\;
    $\MI(G;\, \mathrm{Func}) = \Ent(G) -
    \Ent(G \mid \mathrm{Func})$.
    Sidan $\mathrm{Func}$ er konstant:
    $\Ent(G \mid \mathrm{Func}) = \Ent(G)$.
  \item \textbf{Resultat.}\;
    $\MI(G;\, \mathrm{Func}) = \Ent(G) - \Ent(G) = 0$.
\end{enumerate}

\paragraph{Del~2: $\MI(G;\, M) > 0$.}%
\begin{enumerate}
  \item \textbf{Premiss.}\;
    Frå postulat~\ref{post:L.6.1}:
    $P(\bm{x} \mid m) \neq P(\bm{x})$ for minst eitt
    materiale~$m$.
  \item \textbf{Konsekvens.}\;
    $G$ og $M$ er \emph{ikkje} uavhengige:
    $G \not\perp\!\!\!\perp M$.
  \item \textbf{Fundamental eigenskap.}\;
    $\MI(G;\, M) = 0 \;\iff\; G \perp\!\!\!\perp M$.
  \item \textbf{Resultat.}\;
    Sidan $G \not\perp\!\!\!\perp M$, har vi $\MI(G;\, M) > 0$.
\end{enumerate}

\paragraph{Kombinert:}
$\MI(G;\, M) > 0 = \MI(G;\, \mathrm{Func})$. Materialet ber meir
informasjon om forma enn funksjonen gjer. \qed
\end{proof}


\section{STOLAR-datasettet}

93~stolar frå eit norsk nasjonalmuseum, daterte frå 1280 til 2024,
fordelte på fem materialkategoriar.

\subsection{Geometriske profilar per materiale}

\medskip
\begin{center}
{\footnotesize
\begin{tabular}{@{}lcccc@{}}
\toprule
\textsc{Materiale} & \textsc{Kompaktheit}
  & \textsc{Symmetri} & \textsc{Kurvatur} & \textsc{$n$} \\
\midrule
Mahogni    & 0{,}381 & 0{,}129 & 0{,}105 & 22 \\
Stål       & 0{,}437 & 0{,}118 & 0{,}075 & 19 \\
Kryssfiner & 0{,}455 & 0{,}104 & 0{,}090 & 18 \\
Bøk        & 0{,}445 & 0{,}101 & 0{,}101 & 21 \\
Plast      & 0{,}461 & 0{,}112 & 0{,}121 & 13 \\
\bottomrule
\end{tabular}
}
\end{center}

\medskip
Mønstra er tydelege:
\begin{itemize}[nosep]
  \item \textbf{Mahogni er minst kompakt.} Materialet si høge
    brotstyrke let formgjevaren spreie massen i slankare lemer.
  \item \textbf{Stål har lågast kurvatur.} Flatt stål og I-profil
    tenderer mot det lineære; røyr opnar for kontrollerte kurver.
  \item \textbf{Kryssfiner er mest kompakt blant tremateriala.}
    Laminatet inviterer til skalformer og lukka profilar.
  \item \textbf{Bøk har høgast kurvatur blant tremateriala.}
    Dampbøying opnar for rundare former.
  \item \textbf{Plast har høgast kurvatur totalt.}
    Sprøytestøypt polypropylen opnar for lukka, kompakte geometriar.
    Eames DCW (formkryssfiner, 1946) og Eames Lounge Chair
    (kryssfiner + lær, 1956) demonstrerer korleis substratskiftet
    frå massivt tre til formkryssfiner opna ein ny region i
    formrommet.
\end{itemize}

Effektstorleikane er store (Cohen's $d > 1{,}4$ for dei sterkaste
par). Distribusjonane overlappar knapt.

\subsection{Empiriske informasjonsverdiar}

\begin{align*}
  \MI(G;\, M) &= 0{,}382 \;\text{bit}, \\
  \MI(G;\, \mathrm{Tid}) &= 0{,}168 \;\text{bit}, \\
  \MI(G;\, \mathrm{Geografi}) &= 0{,}079 \;\text{bit}, \\
  \MI(G;\, \mathrm{Func}) &= 0 \;\text{bit (per definisjon).}
\end{align*}
Materialet er den sterkaste prediktoren. Det ber meir enn dobbelt
so mykje informasjon om forma som tida gjer.

\subsection{Materialentropien over tid}

I perioden ca.~1790--1840 utgjorde mahogni 93~prosent av alle
registrerte norskproduserte stolar i nasjonalmuseet
(frå STOLAR-datasettet). Materialentropien kollapsa frå 2{,}5 til
1{,}5~bit. Denne ``mahognitoppen'' illustrerer at ekstremt sterkt
seleksjonstrykk produserer ein ekstremt smal distribusjon.


\section{Designhistorisk døme: Mahognitoppen}

Mellom 1770 og 1840 dominerte mahogni den norske møbelproduksjonen.
Materialet kom frå Karibia og Mellom-Amerika via eit globalt
handelsnettverk forma av kolonialisme, sjøfart og europeisk
prestisjeøkonomi. Merk: stilane Chippendale, Sheraton og
Hepplewhite er primært engelske; dei gjeld for Noreg indirekte,
gjennom import og kopiering.

42~prosent av mahognistolane i det norskproduserte delsettet
kombinerer importert mahogni med lokalt tre: mahogni utanpå, furu
inni (frå STOLAR-datasettet). Stolen er \emph{kolonialt finert}:
globalt prestisjetømmer dekkjer lokal struktur.

Rammeverket fangar denne dynamikken: materialet definerer ein
geometrisk profil, og det kulturelle seleksjonstrykket
konsentrerer populasjonen rundt denne profilen til variansen
kollapsar.


\section{Gjennomarbeidd døme: Rekn ut gjensidig informasjon}

Vi reknar ut $\MI(G;\, M)$ for eit forenkla datasett med
2~materialkategoriar og 2~geometriske kategoriar.

\paragraph{Datasettet.}
20~objekt i funksjonell klasse ``stol'':

\medskip
\begin{center}
{\footnotesize
\begin{tabular}{@{}lcc|c@{}}
\toprule
& \textsc{Kompakt} & \textsc{Open} & \textsc{Sum} \\
\midrule
Tre & 8 & 2 & 10 \\
Stål & 3 & 7 & 10 \\
\midrule
Sum & 11 & 9 & 20 \\
\bottomrule
\end{tabular}
}
\end{center}

\paragraph{Utrekning.}
Marginalfordelingar: $P(\text{Tre}) = P(\text{Stål}) = 0{,}5$;
$P(\text{K}) = 0{,}55$, $P(\text{O}) = 0{,}45$.
\begin{align*}
  \MI(G;\, M)
  &= 0{,}40 \log_2 \frac{0{,}40}{0{,}275}
   + 0{,}10 \log_2 \frac{0{,}10}{0{,}225} \\
  &\quad + 0{,}15 \log_2 \frac{0{,}15}{0{,}275}
   + 0{,}35 \log_2 \frac{0{,}35}{0{,}225} \\
  &= 0{,}216 - 0{,}117 - 0{,}131 + 0{,}223 \\
  &= 0{,}191 \;\text{bit}.
\end{align*}

\paragraph{Tolking.}
$\MI = 0{,}19$~bit er strengt positivt: å vite materialet gjev
informasjon om geometrien. Funksjonen er konstant og gjev
$\MI(G;\, \mathrm{Func}) = 0$.


\section{Øvingar}

\begin{exercise}[Nivå 1]\label{ex:7.1}
Kvifor er $\MI(G;\, \mathrm{Func}) = 0$ innanfor ein funksjonell
klasse? Forklar med eigne ord.
\end{exercise}

\begin{exercise}[Nivå 1]\label{ex:7.2}
Sjå på dei geometriske profilane per materiale. Kva profil ville du
forvente for \emph{aluminium}? Grunngjev svaret.
\end{exercise}

\begin{exercise}[Nivå 2]\label{ex:7.3}
Utvid det gjennomarbeidde dømet: legg til ein tredje material\-kategori
(``plast'') med 10~objekt, der 5~er kompakte og 5~er opne. Rekn
ut den nye $\MI(G;\, M)$. Plast bidreg med null MI sjølv, men
aukar marginalentropien til $M$, noko som senkar det totale
MI-estimatet. Verifiser dette.
\end{exercise}

\begin{exercise}[Nivå 2]\label{ex:7.4}
Materialentropien over tid viser ein kollaps under mahognitoppen.
Formuler dette som eit tilfelle av spreiingsrelasjonen
(Proposisjon~\ref{prop:L.8.4} i neste kapittel).
\end{exercise}

\begin{exercise}[Nivå 3]\label{ex:7.5}
Design eit empirisk eksperiment for å teste postulat~\ref{post:L.6.1}
på ein annan funksjonell klasse enn stolar. Spesifiser geometriske
eigenskapar, materialkategoriar, kontrollvariablar og statistisk
test.
\end{exercise}

\begin{exercise}[Nivå 3]\label{ex:7.6}
Konstruer eit hypotetisk motdøme der
$P(\bm{x} \mid m) = P(\bm{x})$ for alle $m$. Kva ville det krevje
i praksis?
\end{exercise}


%% =====================================================================
%%  KAPITTEL 8: VARIASJONSPRINSIPP
%%  (Tidlegare kap. 7 — flytta hit per PED/W-tilråding)
%% =====================================================================
\chapter{Variasjonsprinsipp}
\label{chap:variasjon}

\begin{quote}\itshape
Etter dette kapittelet skal lesaren kunne:
\begin{itemize}[nosep]
\item Formulere den stiavhengige tilpassingsverdien som eit funksjonale.
\item Utleie Euler--Lagrange-likninga for designtrajektoriar steg for steg.
\item Forklare kvifor stigen modifiserer kva resultat som er moglege.
\item Anvende spreiingsrelasjonen til \aa{} predikere formvarians.
\end{itemize}
\end{quote}

\begin{quote}
\itshape
Det er ikkje posisjonen som tel. Det er stigen.
\end{quote}

No som vi veit at navigasjon er substrat-uavhengig (kap.~6) og at
materialet kanaliserer forma målbart (kap.~7), kan vi spyrje: kva
er den \emph{beste stigen} gjennom formrommet, uansett substrat?


\section{Stigen gjennom formrommet}

Hittil har vi skildra tilpassingslandskapet som ein overflate og
formgjeving som rørsle på denne overflata. Men ein stol er
resultatet av ein prosess --- ein stig gjennom formrommet. Og
prosessen \emph{endrar} kva som er mogleg.


\section{Stiavhengig tilpassing}

\begin{postmark}[Stiavhengig tilpassingsverdi]\label{post:L.8.1}
Tilpassingsverdien til ei realisert form avheng ikkje berre av
posisjonen hennar $\bm{x} \in \MC$, men av heile stigen
$\gamma : [0, T] \to \MC$ som førte dit:
\[
  S[\gamma] \;=\; \int_0^T
    g\!\left(\gamma(t),\; \dot{\gamma}(t),\; t\right) dt
\]
der $g$ er ein lagrangian som avheng av posisjon, hastigheit og tid.
\end{postmark}

\paragraph{Falsifiseringsvilkår.}
Postulatet fell om identiske former med ulik produksjonshistorie
konsekvent har same overlevings- og reproduksjonsrate i
formgjevingsøkologien.


\section{Euler--Lagrange-likninga for designtrajektoriar}

\begin{satmark}[Stasjonær stig]\label{thm:L.8.2}
Dersom $\delta S[\gamma] = 0$, oppfyller $\gamma$ Euler--Lagrange-likninga:
\[
  \frac{\partial g}{\partial \gamma}
  - \frac{d}{dt}\frac{\partial g}{\partial \dot{\gamma}} = 0.
\]
\end{satmark}

\begin{proof}[Bevis, steg for steg]
Vi fylgjer den klassiske variasjonsrekninga. Denne teknikken
generaliserer vanleg optimering til funksjonar av funksjonar
(funksjonalar), og er kanskje ny for lesaren; vi går gjennom
kvart steg eksplisitt.

\paragraph{Steg 1: Definer ein variasjon.}
La $\gamma_\varepsilon(t) = \gamma(t) + \varepsilon \, \eta(t)$
vere ein nabostig, der $\eta(t)$ er ein vilkårleg glatt funksjon
med $\eta(0) = \eta(T) = 0$ (endepunkta er faste).

\paragraph{Steg 2: Formuler stasjonaritetsvilkåret.}
\[
  \frac{d}{d\varepsilon}\bigg|_{\varepsilon=0}
  S[\gamma_\varepsilon] \;=\; 0.
\]

\paragraph{Steg 3: Rekn ut derivaten.}
\begin{align*}
  \frac{d}{d\varepsilon}\bigg|_{0} S[\gamma_\varepsilon]
  &= \int_0^T \!\left(
    \frac{\partial g}{\partial \gamma} \cdot \eta
    + \frac{\partial g}{\partial \dot{\gamma}} \cdot \dot{\eta}
  \right) dt.
\end{align*}

\paragraph{Steg 4: Partiell integrasjon.}
Det andre leddet inneheld $\dot{\eta}$. Vi integrerer delvis:
\begin{align*}
  \int_0^T \frac{\partial g}{\partial \dot{\gamma}}
    \cdot \dot{\eta}\; dt
  &= \left[\frac{\partial g}{\partial \dot{\gamma}}
    \cdot \eta \right]_0^T
    - \int_0^T \frac{d}{dt}
      \frac{\partial g}{\partial \dot{\gamma}}
      \cdot \eta \; dt.
\end{align*}

\paragraph{Steg 5: Grenseleddet forsvinn.}
Sidan $\eta(0) = \eta(T) = 0$, er klammeuttrykket null.

\paragraph{Steg 6: Sett saman.}
\[
  \int_0^T \!\left(
    \frac{\partial g}{\partial \gamma}
    - \frac{d}{dt}\frac{\partial g}{\partial \dot{\gamma}}
  \right) \cdot \eta \; dt \;=\; 0.
\]

\paragraph{Steg 7: Fundamentallemmaet for variasjonsrekninga.}
Sidan $\eta$ er \emph{vilkårleg}, er den einaste måten integralet
kan vere null for alle val av $\eta$ at uttrykket i parentesen er
null overalt (dette er fundamentallemmaet for variasjonsrekninga):
\[
  \frac{\partial g}{\partial \gamma}
  - \frac{d}{dt}\frac{\partial g}{\partial \dot{\gamma}}
  = 0. \qedhere
\]
\end{proof}

\paragraph{Tolking.}
$\frac{\partial g}{\partial \gamma}$ dreg forma mot regionar med
høg tilpassing. $\frac{d}{dt}\frac{\partial g}{\partial
\dot{\gamma}}$ straffar brå endringar. Likninga formaliserer
balansen mellom utforsking og tradisjon.


\section{Stigen modifiserer landskapet}

\begin{satmark}[Stigen endrar kva som er mogleg]\label{thm:L.8.3}
Stiavhengig tilpassing og landskapsminne (frå kap.~4) saman
inneber at å gje form ikkje er å finne den beste \emph{posisjonen},
men den beste \emph{stigen}: prosessen modifiserer kva resultat
som er moglege.
\end{satmark}

\begin{proof}[Grunngjeving]
Frå landskapsminne (kap.~4) veit vi at kvar realisert form
endrar dei framtidige seleksjonstrykka via $\eta_i$:
\[
  s_i(\bm{x},\, t + dt) \;=\;
  s_i(\bm{x},\, t) + \eta_i(\bm{y},\, \bm{x}).
\]
Tilpassingslandskapet ved tidspunkt $T$ er difor eit funksjonale
av \emph{heile} trajektorien.
To ulike trajektoriar som endar i same punkt kan stå overfor ulike
framtidige landskap. Prosessen \emph{modifiserer} kva resultat som
er moglege.
\end{proof}


\section{Spreiingsrelasjonen}

\begin{propmark}[Spreiingsrelasjonen]\label{prop:L.8.4}
La $\sigma^2[\gamma]$ vere variansen i formfordelinga rundt ein
gjeven attraktor i $\MC$, og la $\norm{\bm{s}}$ vere den samla
styrken til seleksjonstrykka ved den attraktoren. Då gjeld:
\[
  \sigma^2 \;\propto\; \frac{1}{\norm{\bm{s}}}.
\]
Sterkare seleksjon produserer smalare fordelingar; svakare
seleksjon produserer breiare.
\end{propmark}

\begin{proof}[Grunngjeving]
Vi modellerer den stasjonære fordelinga nær eit lokalt maksimum
$\bm{x}^*$ som ei Boltzmann-liknande fordeling
(\emph{modelleringsantaknad~1}, ikkje utleidd frå postulata):
\[
  P(\bm{x}) \;\propto\;
  \exp\!\bigl(\beta \, f(\bm{x})\bigr).
\]

\paragraph{Steg 1: Taylor-utvikling.}
Nær $\bm{x}^*$:
\[
  f(\bm{x}) \approx f(\bm{x}^*)
  - \tfrac{1}{2}
  (\bm{x} - \bm{x}^*)^\top H (\bm{x} - \bm{x}^*)
\]
der $H = -\nabla^2 f(\bm{x}^*)$ er positiv definitt.

\paragraph{Steg 2: Sett inn i fordelinga.}
\[
  P(\bm{x}) \propto \exp\!\left(
    -\tfrac{\beta}{2}
    (\bm{x} - \bm{x}^*)^\top H
    (\bm{x} - \bm{x}^*)
  \right).
\]
Dette er ein multivariat normalfordeling med kovariansmatrise
$\Sigma = (\beta H)^{-1}$.

\paragraph{Steg 3: Les av variansen.}
Variansen langs ein vilkårleg akse er
$\sigma_i^2 = (\beta \, \lambda_i)^{-1}$, der $\lambda_i$ er
eigenverdiane til $H$.

\emph{Modelleringsantaknad~2}: Sidan $H$ reflekterer kurvaturen til
dei aggregerte seleksjonstrykka, antek vi at
$\lambda_i \propto \norm{\bm{s}}$. Under denne tilleggsantaknaden
får vi:
\[
  \sigma^2 \;\propto\; \frac{1}{\norm{\bm{s}}}. \qedhere
\]
\end{proof}

\paragraph{Epistemisk status.}
Denne relasjonen kviler på to modelleringsantaknadar
(Boltzmann-fordelinga og proporsjonaliteten
$\lambda_i \propto \norm{\bm{s}}$) som ikkje er utleidde frå
postulata. Ho er ein proposisjon --- ein modellprediksjonen ---
ikkje ein deduktiv konsekvens.


\section{Designhistorisk døme: Formtid og substratskifte}

George Kubler hevda at alt som vert laga er anten ein kopi eller
ein variant av noko som vart laga litt tidlegare
(jf.\ Kublers \emph{prime object}). I vårt rammeverk er dette ein
konsekvens av stiavhengig tilpassing: originalen opna rommet;
kopien fylte det.

Thonet nr.~14 (dampbøygd bøk, 1859) er det sentrale historiske
dømet på korleis teknologisk trykk opnar nye regionar. Eitt
teknologisk gjennombrot --- dampbøying --- endra seleksjonstrykka
permanent. Breuer B3/Wassily-stolen (bøygt stålrøyr, 1925) er
det neste: nytt materiale opna ein ny region, forsterka av
Bauhaus-ideologiens kulturelle trykk. Eames DCW (formkryssfiner,
1946) fullførte triaden: tre som ikkje lenger var avgrensa til
sine tradisjonelle former.


\section{Gjennomarbeidd døme: Spreiingsrelasjonen i praksis}

\paragraph{Kategori A: Militærutstyr (sterkt seleksjonstrykk).}
Ein militærhjelm har strenge krav. Vi estimerer ein høg samla
seleksjonsstyrke $\norm{\bm{s}}_A$.
Prediksjonen: variansen i hjelmgeometri skal vere \emph{låg}.
Observasjonen: militærhjelmar frå ulike produsentar er påfallande
like.

\paragraph{Kategori B: Dekorative møblar (svakt seleksjonstrykk).}
Ein dekorativ lysestake har svake funksjonelle krav.
Prediksjonen: variansen skal vere \emph{høg}.
Observasjonen: lysestakar varierer enormt.

\paragraph{Kvantitativt estimat.}
Med $\beta \norm{\bm{s}}_A = 10$ og
$\beta \norm{\bm{s}}_B = 2$:
\[
  \frac{\sigma^2_B}{\sigma^2_A}
  \;=\; \frac{\norm{\bm{s}}_A}{\norm{\bm{s}}_B}
  \;=\; 5.
\]
Lysestakar har fem gonger høgare geometrisk varians enn
militærhjelmar. Prediksjonen er empirisk testbar.


\section{Øvingar}

\begin{exercise}[Nivå 1]\label{ex:8.1}
Forklar med eigne ord skilnaden mellom å optimere ein
\emph{posisjon} og å optimere ein \emph{stig}.
\end{exercise}

\begin{exercise}[Nivå 1]\label{ex:8.2}
Ein original Wegner-stol og ein fabrikkprodusert replika har
identiske mål. Bruk omgrepet stiavhengig tilpassingsverdi til å
forklare kvifor dei har ulik status.
\end{exercise}

\begin{exercise}[Nivå 2]\label{ex:8.3}
Kva svarar kvart av dei to ledda i Euler--Lagrange-likninga til i
ein konkret designprosess?
\end{exercise}

\begin{exercise}[Nivå 2]\label{ex:8.4}
Vel to funksjonelle klassar --- éin med sterkt seleksjonstrykk og
éin med svakt --- og prediker forholdet mellom variansane deira.
\end{exercise}

\begin{exercise}[Nivå 3]\label{ex:8.5}
Konstruer eit tenkt 2D-formrom med to tradisjonar som startar frå
same punkt men tek ulike ruter. Vis korleis
landskapsminnefunksjonen $\eta_i$ produserer ulike framtidige
landskap.
\end{exercise}

\begin{exercise}[Nivå 3]\label{ex:8.6}
Kublers formsekvens-omgrep svarar til stigen $\gamma$.
Men Kubler antok lineære kjeder. Drøft konsekvensane av at
formrommet er fleirdimensjonalt.
\end{exercise}


%% =====================================================================
%%  KAPITTEL 9: KONSEKVENSAR OG FRAMTID
%% =====================================================================
\chapter{Konsekvensar og framtid}
\label{chap:konsekvensar}

\begin{quote}\itshape
Etter dette kapittelet skal lesaren kunne:
\begin{itemize}[nosep]
\item Bevise at kvar realisert form er eit provisorisk kompromiss.
\item Forklare kvifor tradisjonar med fleire adaptive mekanismar er meir robuste.
\item Vurdere rammeverkets falsifiserbarheit og kva som ville falsifisere det.
\item Samanlikne fire formgjevingspraksisar langs aksane dekning, effektivitet, robustheit og tolkbarheit.
\end{itemize}
\end{quote}

\begin{quote}
\itshape
Ingen form er endeleg.
\end{quote}

Overgangen frå kap.~8 (variasjonsprinsippet) til dette kapittelet
er ein overgang frå analyse til syntese: no brukar vi heile
rammeverket samla.


\section{Provisorisk kompromiss}

\begin{satmark}[Provisorisk kompromiss]\label{thm:L.9.1}
Kvar realisert form er eit \emph{provisorisk} kompromiss:
optimalt for seleksjonstrykka på tidspunktet det vart skapt,
men generisk suboptimalt for alle seinare konfigurasjonar av trykk.
\end{satmark}

\begin{proof}[Grunngjeving]
Frå kompromissteoremet (kap.~2) og postulatet om dynamisk
landskap (kap.~4):
\[
  \grad f(\bm{x}^*(\tau),\, t) \neq \bm{0}
  \quad \text{for } t > \tau \text{ (generisk)}.
\]
Forma som var optimal ved $\tau$ er ikkje lenger optimal ved $t$.
\end{proof}


\section{Varige tradisjonar}

\begin{satmark}[Varige tradisjonar]\label{thm:L.9.2}
Tradisjonar med fleire uavhengige adaptive mekanismar er meir
robuste mot landskapsendringar.
\end{satmark}

\begin{proof}[Grunngjeving]
Ein tradisjon med $k$ uavhengige mekanismar kan respondere på eit
skifte i trykk $s_i$ utan å forstyrre responsen på trykk $s_j$.
La $\Delta k$ vere talet trykk som endrar seg. Andelen av
kapasiteten som vert påverka er $\Delta k / k$: større $k$ gjev
meir redundans. Definer $\Delta k$ som talet av seleksjonstrykk
som endrar seg samstundes. Då er den relative
kapasitetspåverknaden $\Delta k / k$.
\end{proof}


\section{Falsifiserbarheit}

\begin{satmark}[Rammeverket kan felle seg sjølv]%
\label{thm:L.9.3}
Rammeverket er falsifiserbart: dersom eitt einaste
seleksjonstrykk er tilstrekkeleg til å forklare all
observert formvariasjon innanfor ein funksjonell klasse, fell
postulat~2.2 (fleire trykk), og med det fell den deduktive kjeda.
\end{satmark}

\begin{proof}[Den deduktive kjeda]
\begin{enumerate}[nosep]
  \item Fleire trykk $\implies$ variasjon under konstant funksjon.
  \item Fleire trykk + motstridande gradientar $\implies$
    kompromissteoremet $\implies$ fleire toppar.
  \item Dynamisk landskap + kompromiss $\implies$ provisorisk
    kompromiss.
\end{enumerate}
Rammeverket er gyldig \emph{nettopp fordi} det kan feile.
Spreiingsrelasjonen (Proposisjon~\ref{prop:L.8.4}) og
postulat~\ref{post:L.6.1} gjer kvantitative prediksjonar som
kan testast mot data.
\end{proof}


\section{Refleksiv navigasjon}

Alle navigatorane vi har diskutert til no har éin ting til felles:
dei kan ikkje endre sine eigne mål. Menneskelege designarar kan.

Ein handverkar kan spørje seg kvifor ho favoriserer ei viss
stolform, og endre kursen fordi ein refleksjon over eigne mål har
endra kva som tel som ``betre''. Formelt inneber refleksiv
navigasjon at navigatoren endrar aggregeringsfunksjonen $\Phi$ ---
ikkje berre styringsvektoren $\alpha$ --- som ein funksjon av
heile den realiserte banen $\gamma([0, t])$.

Denne kapasiteten er det som gjer designutdanning mogleg: studentane
lærer ikkje berre å navigere betre, men å \emph{endre kva
navigasjon betyr}.


\section{Fire praksisar: Evaluering}

\medskip
\begin{center}
{\footnotesize
\begin{tabular}{@{}lcccc@{}}
\toprule
& \textsc{Handverk} & \textsc{Industri}
  & \textsc{Parametrisk} & \textsc{Gen.\ AI} \\
\midrule
Dekning & Låg & Middels & Middels & Høg \\
Effektivitet & Låg & Middels & Høg & Høg \\
Robustheit & Høg & Middels & Låg & Låg \\
Tolkbarheit & Høg & Middels & Middels & Låg \\
\bottomrule
\end{tabular}
}
\end{center}

\medskip
\begin{description}[nosep]
  \item[Dekning:] Kor stor del av formrommet utforskar navigatoren?
  \item[Effektivitet:] Kor raskt finn navigatoren eit lokalt
    optimum?
  \item[Robustheit:] Kor godt fungerer navigatoren når
    tilpassingslandskapet endrar seg?
  \item[Tolkbarheit:] Kan ein forstå \emph{kvifor} navigatoren
    valde ein posisjon?
\end{description}

\medskip
Ingen navigator er optimal langs alle fire aksane. Å velje
navigator er å velje kva ein prioriterer. Temperaturen i generativ
AI svarar direkte til invers seleksjonsintensitet $\beta$ i
spreiingsrelasjonen: høg temperatur = breiare spreiing.
Diffusjonsmodellar er gradient-ascent i eit approksimert
tilpassingslandskap.


\section{Kva betyr dette for designutdanning?}

\begin{enumerate}
  \item \textbf{Formrommet må ekspliserast.}
    Studentar bør lære å identifisere og kartleggje formrommet.

  \item \textbf{Seleksjonstrykka må namngjevast.}
    Kvar designavgjerd er ei vekting av motstridande trykk.

  \item \textbf{Kompromisset må anerkjennast.}
    Ingen form maksimerer alt samstundes.

  \item \textbf{Substratval er epistemologisk val.}
    Kvart substrat opnar delar av formrommet og stengjer andre.

  \item \textbf{Refleksjon er navigasjonsmodus.}
    Designutdanning som ikkje utviklar refleksiv kapasitet,
    produserer sofistikerte termostatar, ikkje designarar.
\end{enumerate}


\section{Øvingar}

\begin{exercise}[Nivå 1]\label{ex:9.1}
Forklar med eigne ord kva ``provisorisk kompromiss'' betyr. Gje eit
døme på ein artefakt som var optimal på eit tidspunkt men er
suboptimal no.
\end{exercise}

\begin{exercise}[Nivå 1]\label{ex:9.2}
Kvifor er falsifiserbarheit ein styrke, ikkje ein svakheit, for eit
teoretisk rammeverk?
\end{exercise}

\begin{exercise}[Nivå 2]\label{ex:9.3}
Vel éin av dei fire praksisane og analyser ei konkret
designoppgåve (t.d.\ å designe ein lampe) ved å spesifisere
navigator, formrom, seleksjonstrykk, landskap, og eksplorasjon vs.\
eksploitering.
\end{exercise}

\begin{exercise}[Nivå 2]\label{ex:9.4}
Finn eit døme på ein handverkstradisjon som har overlevd i meir enn
200~år, og identifiser minst tre uavhengige adaptive mekanismar.
\end{exercise}

\begin{exercise}[Nivå 3]\label{ex:9.5}
Refleksiv navigasjon inneber at $\Phi$ er ein funksjon av
$\gamma([0, t])$. Kva konsekvensar får det for
Euler--Lagrange-likninga? \emph{Merk: denne øvinga krev kunnskap
utover denne boka.}
\end{exercise}

\begin{exercise}[Nivå 3]\label{ex:9.6}
Temperaturparameteren i generativ AI svarar til $1/\beta$.
Design eit eksperiment med ulike temperaturverdiar og mål den
geometriske variansen. Kva predikerer spreiingsrelasjonen?
\end{exercise}

\begin{exercise}[Nivå 3]\label{ex:9.7}
Vel éi celle i tabellen over fire praksisar (t.d.\
``robustheit for generativ AI = låg'') og konstruer eit argument
for eller mot vurderinga. Kva data treng du?
\end{exercise}


%% =====================================================================
%%  APPENDIKS
%% =====================================================================
\appendix

%% =====================================================================
%%  APPENDIKS A: KORRESPONDANSETABELL
%% =====================================================================
\chapter{Korrespondansetabell}
\label{app:korrespondanse}

Tabellen nedanfor knyter kvart formelt element i denne læreboka
til den tilsvarande posisjonen i \emph{Tractatus Formae} og i
\emph{Formlære: Mathematica}.

\medskip
{\small
\begin{longtable}{@{}p{42mm}p{22mm}p{22mm}@{}}
\toprule
\textsc{Lærebok} & \textsc{Tractatus} & \textsc{Mathematica} \\
\midrule
\endfirsthead
\toprule
\textsc{Lærebok} & \textsc{Tractatus} & \textsc{Mathematica} \\
\midrule
\endhead
\bottomrule
\endfoot

\multicolumn{3}{@{}l}{\itshape Kapittel 1: Formrommet} \\
\addlinespace
Def.\ Formrom & 1.1 & Def.~1.1 \\
Def.\ Regiontypologi & 1.13 & Def.~1.2 \\
Thm.\ Forklaringskrav & 1.14 & Prop.~1.4 \\
\addlinespace

\multicolumn{3}{@{}l}{\itshape Kapittel 2: Seleksjonstrykk} \\
\addlinespace
Def.\ Seleksjonstrykk & 2.1 & Def.~2.1 \\
Def.\ Seleksjonsvektor & 2.14 & Def.~2.2 \\
Def.\ Funksjonell avgr. & 2.4 & Def.~2.3 \\
Thm.\ Var.\ under konst.\ fn. & 2.21 & Thm.~2.5 \\
Post.\ Motstridande grad. & 2.3 & Post.~2.6 \\
Thm.\ Kompromiss & 2.31 & Thm.~2.7 \\
\addlinespace

\multicolumn{3}{@{}l}{\itshape Kapittel 3: Tilpassingslandskapet} \\
\addlinespace
Def.\ Tilpassingsfn. & 3.1 & Def.~3.1 \\
Def.\ Landskapet & 3.2 & Def.~3.2 \\
Thm.\ Fleire haugar & 3.3 & Thm.~3.3 \\
Def.\ Stil & 3.4 & Def.~3.4 \\
\addlinespace

\multicolumn{3}{@{}l}{\itshape Kapittel 4: Dynamikk} \\
\addlinespace
Post.\ Dynamisk landsk. & 4.1 & Post.~4.1 \\
Thm.\ Punktert likevekt & 4.11 & Thm.~4.2 \\
Def.\ Landskapsminne & 4.2 & Def.~4.3 \\
\addlinespace

\multicolumn{3}{@{}l}{\itshape Kapittel 5: Navigatorar} \\
\addlinespace
Def.\ Navigator & 5.1 & Def.~5.1 \\
Def.\ Intelligens & 5.3 & Def.~5.3 \\
Def.\ Kogn.\ lyskjegle & 5.2 & Def.~5.2 \\
\addlinespace

\multicolumn{3}{@{}l}{\itshape Kapittel 6: Substrat og emergens} \\
\addlinespace
Post.\ Substrat-uavh. & 5.12 & Post.~7.1 \\
Posisjonsinvarians & 1.4, 1.41 & Thm.~7.2 \\
Konj.\ Emergent nav. & 5.62 & Conj.~7.3 \\
\addlinespace

\multicolumn{3}{@{}l}{\itshape Kapittel 7: Materialsignaturar} \\
\addlinespace
Post.\ Geometrisk sign. & 8.1 & Post.~6.1 \\
Informasjonsteoremet & 8.11 & Thm.~6.2 \\
STOLAR-data & 8.12--8.18 & Rem.~6.2 \\
\addlinespace

\multicolumn{3}{@{}l}{\itshape Kapittel 8: Variasjonsprinsipp} \\
\addlinespace
Post.\ Stiavh.\ tilp. & 7.1 & Post.~8.1 \\
Stasjonær stig & 7.11 & Thm.~8.2 \\
Stigen endrar landsk. & 7.12 & Thm.~8.3 \\
Prop.\ Spreiingsrel. & 7.2 & Thm.~8.4 \\
\addlinespace

\multicolumn{3}{@{}l}{\itshape Kapittel 9: Konsekvensar} \\
\addlinespace
Provisorisk kompr. & 7.4 & Thm.~9.1 \\
Varige tradisjonar & 7.5 & Thm.~9.2 \\
Falsifiserbarheit & 7.7 & Thm.~9.3 \\
Refleksiv navigasjon & 7.61 & --- (ikkje i Math.) \\
\end{longtable}
}


%% =====================================================================
%%  APPENDIKS B: NOTASJONSOVERSYN
%% =====================================================================
\chapter{Notasjonsoversyn}
\label{app:notasjon}

{\small
\begin{longtable}{@{}p{28mm}p{62mm}@{}}
\toprule
\textsc{Symbol} & \textsc{Tyding} \\
\midrule
\endfirsthead
\toprule
\textsc{Symbol} & \textsc{Tyding} \\
\midrule
\endhead
\bottomrule
\endfoot

$\MC$ & Formrommet (morphospace) for funksjonell klasse~$C$ \\
$\bm{x}(a)$ & Eigenskapsvektoren til objekt~$a$ \\
$B$ & Busett (okkupert) region i~$\MC$ \\
$O$ & Open region: fysisk mogleg men ikkje realisert \\
$F$ & Forboden region: bryt minst éi avgrensing \\
\addlinespace
$s_i$ & Seleksjonstrykk nr.~$i$ \\
$\bm{s}$ & Seleksjonstrykkvektoren $(s_1, \ldots, s_k)$ \\
$k$ & Antal seleksjonstrykk \\
$\norm{\bm{s}}$ & Samla styrke av seleksjonstrykka \\
\addlinespace
$f(\bm{x})$ & Tilpassingsfunksjonen \\
$\Phi$ & Aggregeringsfunksjonen: $\R^k \to \R$ \\
$\bm{w}$ & Vektvektoren for seleksjonstrykk \\
$\FitLand$ & Tilpassingslandskapet \\
\addlinespace
$\gamma$ & Stig (trajektorie) gjennom formrommet \\
$S[\gamma]$ & Stiavhengig tilpassingsfunksjonale \\
$g$ & Lagrangian for designtrajektoriar \\
$\eta_i$ & Påverknadsfunksjon (landskapsminne) \\
\addlinespace
$\Nav$ & Navigator: trippelet $(G, \mu, \alpha)$ \\
$G$ & Målsett (target region) \\
$\mu$ & Avstandsfunksjon \\
$\alpha$ & Styringsfelt (steering field) \\
$\partial \Nav$ & Kognitiv grenseflate (light cone) \\
$I(\Nav)$ & Intelligens: evna til å unngå lokale optima \\
\addlinespace
$\Ent(\cdot)$ & Shannon-entropi \\
$\MI(\cdot\,;\,\cdot)$ & Gjensidig informasjon \\
$\sigma^2$ & Varians i formfordelinga \\
$\beta$ & Seleksjonsintensitet (invers ``temperatur'') \\
$H$ & Hessian-matrisa ($-\nabla^2 f$) \\
$\lambda_i$ & Eigenverdi nr.~$i$ av $H$ \\
\addlinespace
$\Pareto$ & Pareto-fronten \\
$\Rn$ & $n$-dimensjonalt euklidsk rom \\
$\R$ & Dei reelle tala \\
$\grad$ & Gradient ($\nabla$) \\
$\inner{\cdot}{\cdot}$ & Indre produkt \\
$\norm{\cdot}$ & Norm \\
\end{longtable}
}


%% =====================================================================
%%  APPENDIKS C: ØVINGSNØKKEL
%% =====================================================================
\chapter{Øvingsnøkkel}
\label{app:ovingsnokkel}

Denne nøkkelen gjev hint for kvar øving. Fullstendige
utrekningar er utelate med vilje.

\section*{Kapittel 6}

\paragraph{Øving~\ref{ex:6.1}.}
\emph{Hint:} Identifiser $(G, \mu, \alpha)$ for kvart system.
Alle tre oppfyller vilkår (a)--(c) i radikalt ulike substrat.

\paragraph{Øving~\ref{ex:6.2}.}
\emph{Hint:} Prosessen avgjer kva \emph{regionar} som er
tilgjengelege, ikkje kva posisjonen \emph{er}.

\paragraph{Øving~\ref{ex:6.3}.}
\emph{Hint:} $\phi_{jl}$ er kommunikasjonen mellom
teammedlemmane. Det kollektive $G$ emergerer frå forhandlinga.

\paragraph{Øving~\ref{ex:6.4}.}
\emph{Hint:} Vel eigenskapar som fangar topologien: antal hol,
massefyldingsgrad, maks-stress/masse-ratio.

\paragraph{Øving~\ref{ex:6.5}.}
\emph{Hint:} Gje ein algoritme \emph{ubegrensa} reknetid og sjå
om han framleis ikkje kan matche den menneskelege navigatoren.

\paragraph{Øving~\ref{ex:6.6}.}
\emph{Hint:} Gode kandidatar: termittbol, maurtegar,
fuglesvermar, open-source-programvare.


\section*{Kapittel 7}

\paragraph{Øving~\ref{ex:7.1}.}
\emph{Hint:} Ein konstant variabel har null entropi og kan ikkje
gje informasjon om noko anna.

\paragraph{Øving~\ref{ex:7.2}.}
\emph{Hint:} Aluminium (støypt/pressa) gjev verken stål-lineær
eller tre-kurva profil, men ein tredje geometri.

\paragraph{Øving~\ref{ex:7.3}.}
\emph{Hint:} Plast bidreg med null MI sjølv, men aukar
marginalentropien til $M$, noko som senkar det totale MI-estimatet.

\paragraph{Øving~\ref{ex:7.4}.}
\emph{Hint:} Materialpresteige var svært sterkt; $\norm{\bm{s}}$
veks, $\sigma^2$ krympar.

\paragraph{Øving~\ref{ex:7.5}.}
\emph{Hint:} Kontroller for brukstype for å halde funksjonell
klasse konstant. Bruk permutasjonstest.

\paragraph{Øving~\ref{ex:7.6}.}
\emph{Hint:} $P(\bm{x} \mid m) = P(\bm{x})$ krev at alle material
har identiske fysiske eigenskapar --- i praksis umogleg.


\section*{Kapittel 8}

\paragraph{Øving~\ref{ex:8.1}.}
\emph{Hint:} Posisjonsoptimering: ``kva er den beste stolen?''
Stigoptimering: ``kva er den beste \emph{vegen} til stolen?''

\paragraph{Øving~\ref{ex:8.2}.}
\emph{Hint:} Same $\bm{x}$, ulik $\gamma$. Ulik $S[\gamma]$,
ulik verdi.

\paragraph{Øving~\ref{ex:8.3}.}
\emph{Hint:} Fyrste ledd dreg mot høg tilpassing; andre ledd
straffar brå retningsendringar.

\paragraph{Øving~\ref{ex:8.4}.}
\emph{Hint:} Vel t.d.\ medisinsk utstyr (sterkt) og dekorative
vegglamper (svakt). Mål kompaktheit og kurvatur.

\paragraph{Øving~\ref{ex:8.5}.}
\emph{Hint:} Definer $\eta_i$ som ein Gauss-klokke rundt $\bm{y}$.
Vis at ulike ruter forsterkar ulike toppar.

\paragraph{Øving~\ref{ex:8.6}.}
\emph{Hint:} I $n$ dimensjonar kan stigen forgreine seg og sløyfe.
Kublers typologi vert meir kompleks.


\section*{Kapittel 9}

\paragraph{Øving~\ref{ex:9.1}.}
\emph{Hint:} Ein fastlinetelefon var optimal i 1960 men er
suboptimal i 2026.

\paragraph{Øving~\ref{ex:9.2}.}
\emph{Hint:} Eit ufalsifiserbart utsegn forklarer alt og dermed
ingenting.

\paragraph{Øving~\ref{ex:9.3}.}
\emph{Hint:} Følg strukturen frå seksjon om fire praksisar.
Spesifiser kvart element for den valde praksisen.

\paragraph{Øving~\ref{ex:9.4}.}
\emph{Hint:} Kandidatar: japansk sverdsmiing, norsk rosemåling,
italiensk glasblåsing. Identifiser uavhengige mekanismar.

\paragraph{Øving~\ref{ex:9.5}.}
\emph{Hint:} Euler--Lagrange vert ein integrodifferensiallikning.
Denne øvinga krev kunnskap utover boka.

\paragraph{Øving~\ref{ex:9.6}.}
\emph{Hint:} Spreiingsrelasjonen predikerer $\sigma^2 \propto T$
(sidan $T = 1/\beta$). Plott $\sigma^2$ mot $T$.

\paragraph{Øving~\ref{ex:9.7}.}
\emph{Hint:} AI kan ikkje tilpasse seg nye trykk utan re-trening.
Mål kor raskt ein modell vert suboptimal etter eit preferanseskifte.


%% =====================================================================
%%  SLUTTORD
%% =====================================================================

\vfill
\begin{center}
\rule{4cm}{0.4pt}

\medskip
{\small\itshape Denne læreboka er den pedagogiske følgjaren til}\\
{\small\scshape Tractatus Formae}\\
{\small\itshape og}\\
{\small\scshape Formlære: Mathematica}\\[6pt]
{\small\itshape Iver Raknes Finne, 2026}
\end{center}
